\documentclass[10pt,a4paper]{amsart}
\usepackage[margin=19mm]{geometry}
\usepackage{amsmath,amssymb,mathtools}
\usepackage[scr=rsfs,cal=euler]{mathalfa}
\usepackage{xfrac}
\usepackage[unicode,hidelinks,bookmarks=false]{hyperref}
\newcommand{\ie}{i.e.\ }
\newcommand{\cf}{cf.\ }
\newcommand{\resp}{respectively\ }
\newcommand{\Subsection}{\subsection}
\providecommand{\nobreakdash}{\nobreak\hbox{-}}
\setlength{\emergencystretch}{2em}
\multlinegap0pt
\usepackage{amssymb}
\newtheorem{lemmaL}{Lemma}
\newtheorem{corollaryL}{Corollary}
\newtheorem{lemma}{Lemma}
\newtheorem{theorem}{Theorem}
\DeclareMathOperator{\E}{\mathbb{E}}
\title{Power Partitions and Hayman Functions}
\author{Jos\'e L. Fern\'andez, V\'ictor J. Maci\'a}
\date{Source: arXiv:2602.18575v3 (CC BY 4.0)}
\begin{document}
\maketitle

\noindent\emph{Source and licence.} Power Partitions and Hayman Functions. Version arXiv:2602.18575v3 (CC BY 4.0), distributed by arXiv under the Creative Commons Attribution 4.0 International licence, http://creativecommons.org/licenses/by/4.0/, as recorded in the arXiv licence metadata for this exact version and shown on the abstract page https://arxiv.org/abs/2602.18575v3. Full licence notice: \texttt{licenses/arxiv-2602.18575-CC-BY-4.0.txt}.\par\smallskip

\noindent\textit{Editorial scope.} Counted unit: the article's theorem thm:Hayman (source0.tex lines 1074--1077). The article's complete original proof of it is delivered with it (source0.tex lines 1079--1152, one \texttt{proof} environment of 74 lines). No mathematics is written, repaired or re-derived: the statement and the proof are the article's own lines, in the article's own notation, and no retained line is substituted or re-flowed.  Retained uncounted as the source-local context the proof needs: lines 405--410; lines 452--458; lines 488--490; lines 868--878; lines 1037--1044.  Nothing was omitted from any retained span, so the package carries no omission record.  No retained line contains a citation, so the wrapper carries no bibliography and no bibliography entry.  No retained line contains a figure environment, an image-inclusion command or drawing code, so no image is delivered and the package declares no asset dependency.  Editorial formatting only, in addition: an AMS \texttt{amsart} preamble that restates the article's own declarations and macros used by the retained lines, namely the environments \texttt{lemmaL}, \texttt{corollaryL}, \texttt{lemma}, \texttt{theorem} on separate counters, together with the standard packages those lines need.  The retained environments print in this document as Lemmal 1, Corollaryl 1, Lemma 1, Theorem 1 in order of appearance, whereas the article prints the same items with the numbering of its own full document.  The complete native source of the article as distributed in the arXiv e-print for this exact version is bundled unchanged as \texttt{sources/195043/source0.tex}.
\begin{lemmaL}\label{lem:asympt derivatives Fk}
Fix an integer $k\ge 1$. For any integer $m\ge 0$ we have
\[
F_k^{(m)}(-s) \sim \frac{1}{k}\,\zeta(1+1/k)\,\Gamma(m+1/k)\,\frac{1}{s^{m+1/k}}\,, \qquad \text{as } s\downarrow 0\,.
\]
\end{lemmaL}

\begin{equation}\label{eq:mkasympt}
\begin{aligned}
m_k(e^{-s}) &\sim \omega_{k,1}\,\frac{1}{s^{1+1/k}} := \tilde{m}_k(e^{-s}),\\[4pt]
\sigma_k^2(e^{-s}) &\sim \omega_{k,2}\,\frac{1}{s^{2+1/k}} := \tilde{\sigma}_k^2(e^{-s}),
\end{aligned}
\qquad \text{as } s\downarrow 0\,.
\end{equation}

\begin{equation}\label{eq:F3bound}
|F_k'''(-s+\imath \theta)| \le F_k'''(-s), \qquad \text{for any } s>0 \text{ and } \theta\in\mathbb{R}\,,
\end{equation}

\begin{corollaryL}\label{cor:bounds for char Pk}
For any constant $C>0$ there are positive constants $c_1$ and $c_2$ depending
only on~$k$ and~$C$ such that for $s\in(0,\ln 2)$ we have that
\[
\big|\E\!\big(e^{\imath \theta\breve{X}^{[k]}_t}\big)\big| \le
\begin{cases}
e^{-c_1\theta^2}, & \text{if } |\theta|\le C\,\dfrac{1}{s^{1/(2k)}}\,,\\[8pt]
e^{-c_2\frac{1}{s^{1/k}}}, & \text{if } |\theta|\ge C\,\dfrac{1}{s^{1/(2k)}}\,.
\end{cases}
\]
\end{corollaryL}

\begin{lemma}\label{lem:majorarc}
For $f\in\mathcal{K}$ \emph{non-vanishing} in $\mathbb{D}(0,R)$ with Khinchin family $(Z_t)_{t\in[0,R)}$ and
with fulcrum $H(z) = \ln f(e^z)$, for $\Re(z)<\ln R$, and a cut function~$h(t)$, if
\begin{equation}\label{eq:majorarccond}
\lim_{t\uparrow R} \sup_{\varphi\in\mathbb{R}} |H'''(-s+\imath\varphi)| \cdot h(t)^3 = 0\,,
\end{equation}
where $t=e^{-s}$, for $s>\ln(1/R)$, then condition \textup{(H1)} holds.
\end{lemma}

\begin{theorem}\label{thm:Hayman}
The generating function $P_k(z)$ of partitions into $k$-th powers belongs to
the Hayman class.
\end{theorem}

\begin{proof}
We write $t=e^{-s}$ with $s\downarrow 0$ throughout.

Condition (H0) for~$P_k$ holds since from~\eqref{eq:mkasympt}, $\lim_{s\downarrow 0}\sigma_k(e^{-s})=+\infty$.

Next we define an appropriate cut function~$h$. We fix any $\alpha$ in the (nonempty) interval
\begin{equation}\label{eq:alphainterval}
\big(1+1/(3k),\; 1+1/(2k)\big),
\end{equation}
and set $h(e^{-s})=s^{\alpha}$.
\smallskip

We next verify that with this cut~$h$, conditions (H1) and (H2) for the Khinchin family
of~$P_k$ to be in the Hayman class are satisfied. The lower bound  $\alpha > 1+1/(3k)$ is required
to check (H1), while the upper bound  $\alpha < 1+1/(2k)$ is needed to verify (H2).

\medskip

\noindent\textit{Condition \textup{(H1)}.} Equation~\eqref{eq:F3bound} gives
\begin{equation}\label{eq:F3boundbis}
\sup_{\varphi\in\mathbb{R}} \big|F_k'''(-s+\imath\varphi)\big| \le F_k'''(-s), \qquad \text{for } s>0,
\end{equation}
while Lemma~\ref{lem:asympt derivatives Fk} gives that
\[
F_k'''(-s) \sim \omega_{k,3}\, s^{-(3+1/k)}, \qquad \text{as } s\downarrow 0,
\]
and thus since $h(e^{-s})=s^{\alpha}$,
\begin{equation}\label{eq:F3boundtris}
\sup_{\varphi\in\mathbb{R}} \big|F_k'''(-s+\imath\varphi)\big| \cdot h(e^{-s})^3 \le F_k'''(-s)\cdot s^{3\alpha} \sim \omega_{k,3}\, s^{3\alpha-(3+1/k)} \to 0, \qquad \text{as } s\downarrow 0,
\end{equation}
since $\alpha > 1+1/(3k)$. Lemma~\ref{lem:majorarc} then implies that condition (H1) holds for~$P_k$.

\medskip

\noindent\textit{Condition \textup{(H2)}.} The verification of this condition is based upon Corollary~\ref{cor:bounds for char Pk}.

From~\eqref{eq:mkasympt} and the definition of~$h$ we have that
\begin{equation}\label{eq:sigmahasympt}
\begin{aligned}
\sigma_k(e^{-s}) &\sim \frac{\omega_{k,2}^{1/2}}{s^{1+(1/2k)}}\,,\\[4pt]
h(e^{-s})\,\sigma_k(e^{-s}) &\sim \frac{\omega_{k,2}^{1/2}}{s^{1+(1/2k)-\alpha}}\,,
\end{aligned}
\qquad \text{as } s\downarrow 0\,.
\end{equation}

Let $C, c_1, c_2$ be the constants appearing in Corollary~\ref{cor:bounds for char Pk}. Define the following two
sub-arcs of the minor arc
\begin{align*}
I_L(s) &= \big[h(e^{-s})\sigma_k(e^{-s}),\; C s^{-1/(2k)}\big),\\
I_R(s) &= \big[C s^{-1/(2k)},\; \pi\sigma_k(e^{-s})\big).
\end{align*}
For $s$ sufficiently small, these two sub-arcs partition the minor arc, since from~\eqref{eq:sigmahasympt} it
follows that
\begin{equation}\label{eq:arcpartition}
h(e^{-s})\,\sigma_k(e^{-s}) < C s^{-1/(2k)} < \pi\,\sigma_k(e^{-s}), \qquad \text{for  } s>0 \text{ small enough}.
\end{equation}

For $|\theta|\in I_L(s)$, since $|\theta|\ge h(e^{-s})\sigma_k(e^{-s})$, Corollary~\ref{cor:bounds for char Pk} gives that
\begin{equation}\tag{$\flat_L$}
|\E(e^{\imath\theta\breve{X}^{[k]}_t})| \le e^{-c_1(h(e^{-s})\sigma_k(e^{-s}))^2}.
\end{equation}

For $|\theta|\in I_R(s)$, Corollary~\ref{cor:bounds for char Pk} gives directly that
\begin{equation}\tag{$\flat_R$}
|\E(e^{\imath\theta\breve{X}^{[k]}_t})| \le e^{-c_2/s^{1/k}}.
\end{equation}

Let $\beta = \min\{2+1/k-2\alpha,\; 1/k\} > 0$. Appealing to~\eqref{eq:sigmahasympt}, we may combine $(\flat_L)$ and
$(\flat_R)$ and obtain, for a certain constant $c_3>0$ and for $s$ small enough, that
\[
\sup_{h(t)\sigma_k(t)<|\theta|\le\pi\sigma_k(t)} \big|\E\!\big(e^{\imath\theta\breve{X}_t}\big)\big| \le e^{-c_3/s^{\beta}}.
\]
This bound combined with the asymptotics for~$\sigma_k$ of~\eqref{eq:sigmahasympt} gives that (H2) holds.
\end{proof}

\end{document}
